\documentclass[11pt,a4paper]{article}
\usepackage[utf8]{inputenc}
\usepackage[T1]{fontenc}
\usepackage[french]{babel}
\frenchbsetup{StandardLists=true}
\usepackage[left=2cm,right=2cm,top=2cm,bottom=2cm]{geometry}
\usepackage[dvipsnames]{xcolor}
\usepackage{fouriernc}
\usepackage{amsmath,amssymb,amsfonts}
\usepackage{array,tabularx,booktabs,multirow}
\usepackage{colortbl}
\usepackage{graphicx}
\usepackage{mdframed}
\usepackage{enumitem}
\usepackage{fancyhdr}
\usepackage{chemformula}
\usepackage{awesomebox}
\setlength{\aweboxvskip}{0mm}
\usepackage{tcolorbox}
\tcbuselibrary{skins}
\usepackage{fontawesome5}

\definecolor{exoheader}{HTML}{1E5AA0}
\definecolor{exolight}{HTML}{DCE8FF}
\definecolor{warnorange}{RGB}{220,100,0}
\definecolor{problemebg}{HTML}{FFFDE7}

\renewcommand{\arraystretch}{1.8}
\setlength{\extrarowheight}{1mm}

\pagestyle{fancy}
\renewcommand\headrulewidth{1pt}
\fancyhead[L]{Académie de Besançon}
\fancyhead[C]{\textbf{TP --- Titrage conductimétrique du NaCl}}
\fancyhead[R]{Terminale / BTS Chimie}
\fancyfoot[C]{page \thepage}
\fancyfoot[R]{2025 -- 2026}
\fancyfoot[L]{Alexis RUIZ}

\newcommand{\sectitle}[1]{%
  \vspace{6pt}
  \noindent\colorbox{exoheader}{\parbox{\dimexpr\linewidth-2\fboxsep}%
    {\color{white}\bfseries\large\strut\quad #1}}
  \vspace{4pt}
}
\newcommand{\subsectitle}[1]{%
  \vspace{4pt}
  \noindent\colorbox{exolight}{\parbox{\dimexpr\linewidth-2\fboxsep}%
    {\color{exoheader}\bfseries\strut\quad #1}}
  \vspace{2pt}
}
\newcounter{question}
\newcommand{\question}{%
  \stepcounter{question}%
  \noindent\textcolor{exoheader}{\textbf{Question \thequestion.}}\quad
}

\begin{document}

\noindent\textbf{NOM :}\hfill\textbf{Prénom :}\hfill\textbf{Classe :}\hfill\phantom{x}\\[0.2cm]
\hrule\vspace{0.2cm}

\begin{center}
  {\large\bfseries\textcolor{exoheader}{Travaux Pratiques de Chimie}}\\[4pt]
  {\LARGE\bfseries Titrage conductimétrique du sérum physiologique}\\[2pt]
  {\small Académie de Besançon --- Belfort}
\end{center}
\hrule\vspace{6pt}

\begin{tcolorbox}[enhanced, colback=problemebg, colframe=exoheader, boxrule=1.5pt,
  title={\faFlask\quad Problématique},
  fonttitle=\bfseries\color{white},
  attach boxed title to top left={yshift=-2mm, xshift=4mm},
  boxed title style={colback=exoheader, colframe=exoheader}]
\textit{Le sérum physiologique est une solution de \ch{NaCl} à 9~g/L (0,9\% en masse).
\textbf{Peut-on vérifier cette teneur par titrage conductimétrique par précipitation avec \ch{AgNO3} ?}}
\end{tcolorbox}

\vspace{6pt}

% ══════════════════════════════════════════════════════════════════════════════
\sectitle{Partie I --- Rappels théoriques}
% ══════════════════════════════════════════════════════════════════════════════

\subsectitle{1. Principe du titrage conductimétrique}

\begin{mdframed}[backgroundcolor=exolight, linecolor=exoheader, linewidth=1.2pt]
On titre les ions \ch{Cl-} par une solution de nitrate d'argent \ch{AgNO3} :
\[
  \ch{Ag+_{(aq)}} + \ch{Cl-_{(aq)}} \longrightarrow \ch{AgCl_{(s)}} \quad \text{(précipité blanc)}
\]
La conductivité $\sigma$ varie en fonction du volume versé, avec un changement de pente
au \textbf{point d'équivalence}.
\end{mdframed}

\vspace{4pt}
\question Expliquer l'évolution de la conductivité avant et après l'équivalence
(utiliser les conductivités molaires ioniques en annexe). \dotfill

\vspace{0.3cm}\noindent\dotfill

\vspace{0.3cm}\noindent\dotfill

\question Quel est l'intérêt du titrage conductimétrique par rapport à un titrage avec indicateur coloré ? \dotfill

\vspace{0.3cm}\noindent\dotfill

% ══════════════════════════════════════════════════════════════════════════════
\sectitle{Partie II --- Protocole expérimental}
% ══════════════════════════════════════════════════════════════════════════════

\subsectitle{1. Préparation de la solution de sérum diluée}

\begin{enumerate}[leftmargin=1.5cm, label=\textcolor{exoheader}{\textbf{\arabic*.}}]
  \item Prélever $V_{\text{sérum}} = 10{,}0$~mL de sérum physiologique à la pipette jaugée.
  \item Transvaser dans une fiole jaugée de 100~mL, compléter à la jauge avec de l'eau distillée.
  \item Boucher et retourner 3 fois. Verser dans le bécher étiqueté \og sérum dilué \fg{}.
\end{enumerate}

\noindent\textcolor{exoheader}{\faCheckSquare\quad \textbf{Appel n°1} : Faire vérifier la solution diluée par le professeur.}

\vspace{4pt}
\subsectitle{2. Mise en place du montage}

\begin{enumerate}[leftmargin=1.5cm, label=\textcolor{exoheader}{\textbf{\arabic*.}}]
  \item Rincer la burette avec \ch{AgNO3} ($C_b = 0{,}0200$~mol/L), puis la remplir à 0,0~mL.
  \item Placer $V_a = 20{,}0$~mL de sérum dilué dans le bécher (pipette jaugée).
  \item Ajouter 30~mL d'eau distillée (pour immerger la sonde).
  \item Plonger la sonde conductimétrique et le barreau aimanté ; mettre l'agitation en marche.
\end{enumerate}

\subsectitle{3. Titrage conductimétrique}

\begin{enumerate}[leftmargin=1.5cm, label=\textcolor{exoheader}{\textbf{\arabic*.}}]
  \item Mesurer et noter la conductivité initiale ($V = 0$~mL).
  \item Ajouter \ch{AgNO3} par tranches de 1,0~mL ; attendre la stabilisation ; noter $\sigma$.
  \item Autour du changement de pente, serrer les ajouts à 0,5~mL.
  \item Continuer jusqu'à $V = 25$~mL ou nette augmentation de la conductivité.
\end{enumerate}

\noindent\textcolor{exoheader}{\faCheckSquare\quad \textbf{Appel n°2} : Faire vérifier la courbe par le professeur.}

% ══════════════════════════════════════════════════════════════════════════════
\newpage
\sectitle{Partie III --- Tableau de mesures et tracé de la courbe}
% ══════════════════════════════════════════════════════════════════════════════

\begin{center}
\begin{tabularx}{\linewidth}{|>{\centering\arraybackslash\columncolor{exolight}}X
                              |>{\centering\arraybackslash}X
                              |>{\centering\arraybackslash\columncolor{exolight}}X
                              |>{\centering\arraybackslash}X
                              |>{\centering\arraybackslash\columncolor{exolight}}X
                              |>{\centering\arraybackslash}X|}
\hline
\rowcolor{exoheader}
\color{white}$V$ (mL) & \color{white}$\sigma$ (mS/cm) &
\color{white}$V$ (mL) & \color{white}$\sigma$ (mS/cm) &
\color{white}$V$ (mL) & \color{white}$\sigma$ (mS/cm) \\
\hline
0,0 & & 9,0 & & 16,0 & \\
\hline
1,0 & & 10,0 & & 17,0 & \\
\hline
2,0 & & 11,0 & & 18,0 & \\
\hline
3,0 & & 12,0 & & 19,0 & \\
\hline
4,0 & & 13,0 & & 20,0 & \\
\hline
5,0 & & 14,0 & & 21,0 & \\
\hline
6,0 & & 14,5 & & 22,0 & \\
\hline
7,0 & & 15,0 & & 23,0 & \\
\hline
8,0 & & 15,5 & & 25,0 & \\
\hline
\end{tabularx}
\end{center}

\vspace{4pt}
\question Tracer la courbe $\sigma = f(V)$ ci-dessous (ou sur papier millimétré).

\vspace{5cm}
\begin{center}
\textit{(graphique $\sigma = f(V)$)}
\end{center}

\noindent\textcolor{exoheader}{\faCheckSquare\quad \textbf{Appel n°3} : Faire vérifier le tracé par le professeur.}

% ══════════════════════════════════════════════════════════════════════════════
\sectitle{Partie IV --- Exploitation des résultats}
% ══════════════════════════════════════════════════════════════════════════════

\subsectitle{1. Lecture graphique du point d'équivalence}

\question Déterminer graphiquement le volume équivalent $V_{éq}$.

\[
  V_{éq} = \dotfill \text{ mL}
\]

\subsectitle{2. Calcul de la concentration en NaCl}

\question À l'équivalence : $C_b \times V_{éq} = C_a \times V_a$.
Calculer la concentration molaire $C_a$ du sérum dilué.
($C_b = 0{,}0200$~mol/L ; $V_a = 20{,}0$~mL)

\vspace{1.5cm}

\[
  C_a = \dotfill \text{ mol/L}
\]

\subsectitle{3. Concentration du sérum commercial}

\question Le sérum a été dilué 10 fois. Calculer $C_{\text{sérum}}$ puis la concentration massique.
($M_{\ch{NaCl}} = 58{,}5$~g/mol)

\vspace{1.5cm}

\[
  C_m = \dotfill \text{ g/L}
\]

\question Comparer avec la valeur attendue (9~g/L) et conclure. \dotfill

\vspace{0.3cm}\noindent\dotfill

% ══════════════════════════════════════════════════════════════════════════════
\newpage
\sectitle{Annexe --- Données}
% ══════════════════════════════════════════════════════════════════════════════

\begin{mdframed}[backgroundcolor=exolight, linecolor=exoheader, linewidth=1.2pt]
\textbf{Conductivités molaires ioniques} à 25°C :
\begin{itemize}
  \item \ch{Cl-} : $\lambda° = 76{,}4$~mS$\cdot$m²/mol
  \item \ch{NO3-} : $\lambda° = 71{,}5$~mS$\cdot$m²/mol
  \item \ch{Ag+} : $\lambda° = 61{,}9$~mS$\cdot$m²/mol
  \item \ch{Na+} : $\lambda° = 50{,}1$~mS$\cdot$m²/mol
\end{itemize}
\textbf{Masses molaires :} $M_{\ch{NaCl}} = 58{,}5$~g/mol ; $M_{\ch{AgNO3}} = 169{,}9$~g/mol

\textbf{Réactif titrant :} Solution \ch{AgNO3}, $C_b = 0{,}0200$~mol/L

\textbf{Réaction de titrage :} \ch{Ag+_{(aq)}} + \ch{Cl-_{(aq)}} $\to$ \ch{AgCl_{(s)}}
\end{mdframed}

\end{document}
